\section{Elimination certificates that produce complete affine rational-point lists}
\label{sec:elimination}

The preceding proof-transport discipline becomes particularly useful when
a classical finiteness argument has supplied auxiliary equations.  Existing
packets already contain complete rational-root scans
\cite{repo-geometry}, determinant norms for monic plane equations
\cite{repo-demand}, and sharp weighted refinements for several curve
families \cite{repo-coherent}.  We combine those finite mechanisms in a
different way: an arbitrary plane presentation is put into a suitable
coordinate system by a bounded search, and its elimination polynomial
produces a \emph{complete list} of rational intersection points.  A uniform
coefficient envelope on auxiliary equations then suffices to produce a
complete list for the target set, without choosing an auxiliary equation
for each point.  This makes explicit a stronger interface than the
finite-tuple covering premise in the Siegel packet \cite{repo-siegel}.

The algebra is standard elimination theory in origin
\cite{sturmfels}.  The contribution here is the full arithmetic
construction, its exact guards, its quantifier-sensitive composition, and
its induction audit.  No general Siegel or Faltings theorem is asserted.

\subsection{Exact inputs, guards, and bounded coordinates}

A bivariate rational polynomial is a finite coefficient array
\(F=(f_{ij})\), with default value zero outside the encoded array.  Its
value is the finite sum
\[
 F(x,y)=\sum_{i,j}f_{ij}x^iy^j.
\]
Rationals have canonical signed numerators and positive denominators.
Equality, order, polynomial evaluation, and every array operation used
below are primitive recursive.  The total degree of a nonzero polynomial
is the largest \(i+j\) having \(f_{ij}\ne0\); a separate Boolean flag
recognizes the zero polynomial.  All degrees below are total degrees
unless the variable is displayed explicitly.

The final algorithm is total and returns either a successful list or an
\emph{unsupported} flag.  If either polynomial is a nonzero constant, it
returns the empty list successfully.  If one is the zero polynomial and
neither is a nonzero constant, it returns unsupported.  Thus its main
branch has nonzero polynomials \(F,G\) of degrees \(d,e\ge1\).
The unsupported flag is never interpreted as an empty solution set.

For an integer \(c\ge0\), make the shear
\[
 U=X+cY,\qquad V=Y,\qquad
 F_c(U,V)=F(U-cV,V),\quad G_c(U,V)=G(U-cV,V).
 \tag{E1}\label{eq:elim-shear}
\]
The inverse map is \((U,V)\mapsto(U-cV,V)\); both maps preserve rational
coordinates and are mutually inverse by direct substitution.  Binomial
expansion computes the coefficient arrays.

\begin{lemma}[A bounded generic direction]\label{lem:elim-shear}
Let \(F_d,G_e\) be the nonzero highest homogeneous parts of \(F,G\).
There is a \(c\in\{0,\ldots,d+e\}\) such that
\[
 a(c):=F_d(-c,1)\ne0,\qquad b(c):=G_e(-c,1)\ne0.
 \tag{E2}\label{eq:elim-leading}
\]
The first such \(c\) is a primitive recursive function of the coefficient
arrays.  For that choice,
\[
 F_c=\sum_{i=0}^d f_i(U)V^i,\quad
 G_c=\sum_{j=0}^e g_j(U)V^j,
\]
where
\[
 f_d=a(c),\quad g_e=b(c),\quad
 \deg_U f_i\le d-i,\quad \deg_U g_j\le e-j
 \tag{E3}\label{eq:elim-degree-rows}
\]
for nonzero coefficients.  In particular, every rational specialization
\(F_c(u,V)\) has degree exactly \(d\).
\end{lemma}

\begin{proof}
The coefficients of the polynomial \(a(T)=F_d(-T,1)\) are the coefficients
of \(F_d\), with their signs changed according to the exponent of \(X\).
Thus \(a\) is nonzero and has degree at most \(d\).  Similarly \(b\) is
nonzero of degree at most \(e\).  Their product is nonzero: the product of
their leading coefficients is nonzero in \(\Q\).

We recall the finite root bound in the form needed here.  If a nonzero
degree-\(r\) rational polynomial vanishes at distinct supplied numbers
\(t_1,\ldots,t_s\), synthetic division gives
\(P(T)=(T-t_1)P_1(T)\).  Evaluating at \(t_j\ne t_1\) shows
\(P_1(t_j)=0\), because \(t_j-t_1\ne0\).  Continue along the finite
list.  A nonzero polynomial cannot lose more than its degree in these
divisions, so \(s\le r\).  Each quotient and division identity is finite
coefficient data; no search for an unspecified root is involved.

Consequently \(a(T)b(T)\) cannot vanish at all \(d+e+1\) candidates
\(0,\ldots,d+e\).  A scan with a success flag, frozen at the first passing
candidate, returns the required \(c\).  Finally the shear preserves the
bound on total degree.  A term containing \(V^d\) in \(F_c\) cannot
contain a positive power of \(U\), and its coefficient is
\(F_d(-c,1)\).  Every coefficient of \(V^i\) has \(U\)-degree at most
\(d-i\).  This proves all the assertions.
\end{proof}

Henceforth \(c\) always denotes this canonical first passing value.  The
constant leading coefficient is useful twice: it controls the elimination
degree, and it prevents the polynomial in the second coordinate from
vanishing identically at an exceptional value of the first coordinate.

\subsection{The finite elimination certificate and its degree}

Put \(N=d+e\).  Let \(S(U)\) be the \(N\)-by-\(N\) matrix whose rows
are coefficient vectors, in columns \(V^0,\ldots,V^{N-1}\), of
\[
 F_c,VF_c,\ldots,V^{e-1}F_c,
 \qquad G_c,VG_c,\ldots,V^{d-1}G_c.
\]
Define the elimination polynomial
\[
 R(U)=\det S(U).
 \tag{E4}\label{eq:elim-resultant}
\]
This is a Sylvester determinant; other conventional row orders change at
most its sign.  We use the permutation-sum definition, so the certificate
can be checked by finite rational polynomial arithmetic.  In the main
branch, \(\Good(F,G)\) means that the coefficient array of \(R\) is
nonzero.  It is a decidable primitive recursive test.  It is not presumed
to succeed on every pair.

\begin{lemma}[The total-degree audit]\label{lem:elim-degree}
If \(R\ne0\), then \(\deg R\le de\).
\end{lemma}

\begin{proof}
Consider one nonzero term of the determinant expansion.  In the row
\(V^rF_c\), an entry selected from column \(k\) is \(f_{k-r}\), and
its degree is at most \(d-(k-r)\).  In the row \(V^sG_c\), the
corresponding bound is \(e-(k-s)\).  A permutation uses each column
\(0,\ldots,d+e-1\) once.  Therefore the degree of the selected product
is at most
\begin{align*}
 2de-\sum_{k=0}^{d+e-1}k
       +\sum_{r=0}^{e-1}r+\sum_{s=0}^{d-1}s
 &=2de-\frac{(d+e)(d+e-1)}2
          +\frac{e(e-1)+d(d-1)}2\\
 &=de.
\end{align*}
A zero selected entry makes the determinant term zero and requires no
degree calculation.  Finite addition cannot create a coefficient above
the maximum support bound.  This proves the claim without any assumption
that leading determinant terms do not cancel.
\end{proof}

\begin{lemma}[A common point gives an elimination root]
\label{lem:elim-root}
If \(F(x,y)=G(x,y)=0\), then \(R(x+cy)=0\).
\end{lemma}

\begin{proof}
Write \(u=x+cy\), \(v=y\).  The column
\(w=(1,v,\ldots,v^{N-1})^{\mathsf T}\) satisfies \(S(u)w=0\): each
row evaluates a shifted copy of \(F_c\) or \(G_c\) at \(v\).
Multiply by the adjugate of \(S(u)\).  The finite adjugate identity gives
\(R(u)w=0\), and its first component is \(R(u)=0\).
Determinant specialization and the adjugate identity follow from their
finite polynomial expansions, so no assertion about a completed
algebraic closure is used.
\end{proof}

For complete enumeration this is already enough.  For the sharp
cardinality bound one must account for several points with the same
\(u\)-coordinate.  Counting only distinct roots of \(R\) would omit
that information.

\begin{lemma}[A fibre contributes its full multiplicity]
\label{lem:elim-fibre}
If \(v_1,\ldots,v_r\) are distinct rationals and
\(F_c(u,v_j)=G_c(u,v_j)=0\) for every \(j\), then
\[
 (U-u)^r\mid R(U).
 \tag{E5}\label{eq:elim-fibre-power}
\]
\end{lemma}

\begin{proof}
The polynomial \(F_c(u,V)\) has degree exactly \(d\), so the finite
root bound gives \(r\le d\le N\).  Form a constant \(N\)-by-\(N\)
matrix \(W\).  Its first \(r\) columns are
\[
 (1,v_j,\ldots,v_j^{N-1})^{\mathsf T},\qquad 1\le j\le r;
\]
its remaining columns are the standard basis columns
\(e_r,\ldots,e_{N-1}\), numbering coordinates from zero.  The
upper-left block is a Vandermonde matrix, the upper-right block is zero,
and the lower-right block is the identity.  Thus
\[
 \det W=\prod_{i<j}(v_j-v_i)\ne0.
\]
The Vandermonde formula is itself a finite polynomial identity: as a
polynomial in the last ordinate its determinant has the previous
ordinates as roots and the preceding Vandermonde determinant as its
leading coefficient; induction proves the displayed product.

Each of the first \(r\) columns of \(S(U)W\) vanishes entrywise at
\(U=u\).  Synthetic division therefore factors \(U-u\) from every
entry in each such column.  Factoring a scalar polynomial from these
\(r\) columns in the permutation-sum determinant proves
\((U-u)^r\mid\det(S(U)W)\).  Multiplicativity of the finite determinant
gives \(\det(S(U)W)=R(U)\det W\).  Dividing the quotient polynomial by
the nonzero rational \(\det W\) proves \eqref{eq:elim-fibre-power}.
\end{proof}

\begin{theorem}[Finite arithmetic B\'ezout bound]\label{thm:elim-bezout}
Assume \(\Good(F,G)\).  Every finite list of pairwise distinct rational
common zeros of \(F,G\) has length at most \(de\).
\end{theorem}

\begin{proof}
Apply the shear to the list and group entries by their first coordinate,
using decidable equality of canonical rationals.  Let the distinct first
coordinates be \(u_1,\ldots,u_t\), with group sizes
\(r_1,\ldots,r_t\).  Within a group the second coordinates are distinct,
because the shear is injective.  Lemma~\ref{lem:elim-fibre} gives
\((U-u_j)^{r_j}\mid R\) for every \(j\).

For completeness, coprimeness of these powers has explicit finite
witnesses.  If \(a\ne b\), put
\[
 L=\frac{U-a}{b-a},\qquad M=-\frac{U-b}{b-a},\qquad L+M=1.
\]
Expand \(1=(L+M)^{r+s-1}\).  A term whose exponent of \(L\) is at
least \(r\) is divisible by \((U-a)^r\); every other term has exponent
of \(M\) at least \(s\) and is divisible by \((U-b)^s\).  Grouping
the terms gives polynomial B\'ezout coefficients for the two powers.
Multiplying such identities supplies coprime cancellation for a finite
product.  Hence
\[
 \prod_{j=1}^t(U-u_j)^{r_j}\mid R(U).
\]
The product is monic of degree \(\sum_jr_j\).  Since \(R\ne0\),
degree comparison and Lemma~\ref{lem:elim-degree} yield
\[
 \operatorname{length}(\text{list})
     =\sum_jr_j\le\deg R\le de.
\]
All factorization witnesses and grouped lists are finite coded data.
\end{proof}

The statement needs neither an irreducibility oracle nor an unproved
claim about common components.  In classical algebra, absence of a common
nonconstant factor implies the good-shear resultant is nonzero.  That
characterization is unnecessary here: \(R\ne0\) is checked directly.
When it fails, even a finite or empty rational intersection may be
unsupported.  Completeness is asserted on the explicitly recognized
successful domain, not for all pairs of equations.

\subsection{From elimination to a complete list and height bound}

We include the rational-root argument to make the resulting positive
finiteness theorem self-contained.  This ingredient is already available
in the repository \cite{repo-geometry}; the new step is its composition
with the uniformly chosen two-coordinate elimination.

\begin{lemma}[Complete bounded rational-root scan]
\label{lem:elim-rational-roots}
For each nonzero rational polynomial \(P(T)\), a fixed primitive
recursive program returns exactly its rational roots, without duplicates.
\end{lemma}

\begin{proof}
Multiply all coefficients by their positive denominator product.  This
gives a nonzero integer polynomial with the same roots.  By bounded scans
find its least and greatest nonzero coefficient indices \(k,n\).  If
\(k>0\), record the root zero.  Write
\[
 P(T)=T^kQ(T),\qquad
 Q(T)=a_0+a_1T+\cdots+a_mT^m,
 \quad a_0a_m\ne0,quad m=n-k.
\]
If \(m=0\), there are no further roots.  Otherwise take a nonzero root
\(p/q\) in lowest terms, with \(q>0\).  Multiplication by \(q^m\)
gives the integer identity
\[
 a_mp^m+a_{m-1}p^{m-1}q+\cdots+a_0q^m=0.
 \tag{E6}\label{eq:elim-root-identity}
\]
It follows that \(p\mid a_0q^m\) and \(q\mid a_mp^m\).
The Euclidean algorithm gives a B\'ezout equation for the coprime
\(p,q\).  If an integer is coprime to two integers separately, multiplying
their B\'ezout equations proves it coprime to their product; induction
therefore proves coprimeness with the required powers.  Coprime
cancellation in \eqref{eq:elim-root-identity} now gives
\[
 p\mid a_0,\qquad q\mid a_m,
 \qquad 0<|p|\le|a_0|,quad 1\le q\le|a_m|.
 \tag{E7}\label{eq:elim-root-box}
\]
Scan this finite integer rectangle, retain coprime pairs whose exact
substitution into \(P\) is zero, and canonicalize.  The final equation
test proves soundness, and \eqref{eq:elim-root-box} proves that every
nonzero rational root is tested.  The separately recorded zero is a root
exactly when \(k>0\).  Canonical deduplication preserves membership.

A nonzero constant returns the empty list.  The zero polynomial has no
finite root-list contract and receives a failure flag in the total
wrapper.  Every scan bound is a numerical function of already computed
coefficients, so all operations are primitive recursive.
\end{proof}

\begin{theorem}[A positive plane-intersection theorem]
\label{thm:elim-complete}
There is a fixed primitive recursive function \(\operatorname{Intersect}\)
such that, if \(\Good(F,G)\), its output is a duplicate-free list \(L\)
with
\[
 (x,y)\in L\quad\Longleftrightarrow\quad F(x,y)=G(x,y)=0
 \quad\text{for all }x,y\in\Q.
 \tag{E8}\label{eq:elim-complete}
\]
Moreover \(|L|\le de\).  It also computes an upper bound on the canonical
codes, and on the naive coordinate heights, of all these points.
\end{theorem}

\begin{proof}
The algorithm is explicit:
\begin{enumerate}
\item Compute the canonical shear \(c\) and its nonzero polynomial \(R\).
\item Enumerate all rational roots \(u\) of \(R\), using
Lemma~\ref{lem:elim-rational-roots}.
\item For each \(u\), enumerate all rational roots \(v\) of
\(F_c(u,V)\).  Its leading coefficient is the fixed nonzero rational
\(a(c)\), so this call never encounters the zero polynomial.
\item Retain the pairs satisfying \(G_c(u,v)=0\), replace them by
\((u-cv,v)\), and remove duplicates.
\end{enumerate}
Every retained point satisfies both original equations by the inverse
shear.  Conversely let \(F(x,y)=G(x,y)=0\).  Put \(u=x+cy\) and
\(v=y\).  Lemma~\ref{lem:elim-root} puts \(u\) in the first complete
root list.  The equation \(F_c(u,v)=0\) puts \(v\) in that fibre's
complete list.  The other equation passes the final filter, and the
inverse shear returns \((x,y)\).  This proves
\eqref{eq:elim-complete}.  Theorem~\ref{thm:elim-bezout} applies to the
duplicate-free output and gives its length bound.

For a canonical rational \(z=p/q\), put \(h(z)=\max(|p|,q)\), and
for a pair put \(h(x,y)=\max(h(x),h(y))\).  The maximum of these heights
over \(L\), with zero for an empty list, bounds every rational common
point by completeness.  Taking the maximum of the point codes works in
the same way.  Neither maximum is an unbounded search.

The outer root rectangle is determined by \(R\); each inner rectangle
is determined by the coefficients of \(F_c(u,V)\) already computed for
that outer entry.  Nesting loops with such primitive recursive bounds
again gives a primitive recursive program.  Thus the conclusion includes
an actual list and height producer, rather than only a cardinality bound.
\end{proof}

The theorem is a zero-dimensional result: successful elimination has
proved that the relevant intersection has a finite rational list.  It
does not bound all rational points on a curve in the absence of auxiliary
equations.

\subsection{Two cover interfaces and a bounded-envelope theorem}

Fix a nonzero nonconstant \(F\), and a fixed primitive recursive
decidable predicate \(A(x,y)\) on rational codes.  Ordinary numerical
parameters of this predicate may be carried through every program and
proof below; they are suppressed in the notation.  Integrality of both
coordinates is one example.  Define the target set by
\[
 S(x,y)\quad:\Longleftrightarrow\quad F(x,y)=0\ \wedge\ A(x,y).
\]
This is notation for a numerical predicate, not an additional set sort.

\begin{corollary}[A supplied global auxiliary cover]
\label{cor:elim-global-cover}
Suppose a finite list \(G_1,\ldots,G_m\) is supplied, each pair
\((F,G_i)\) passes the good-resultant test, and the following mathematical
premise has been proved:
\[
 \forall x,y\in\Q\quad
 S(x,y)\ \longrightarrow\ \bigvee_{i=1}^mG_i(x,y)=0.
 \tag{E9}\label{eq:elim-global-cover}
\]
Then a primitive recursive function of these data produces a complete
list of the target points and a height bound.  Its length is at most
\(\sum_i d\deg G_i\).
\end{corollary}

\begin{proof}
Concatenate the lists \(\operatorname{Intersect}(F,G_i)\), filter by
\(A\), and deduplicate.  A retained point satisfies \(F=0\) and \(A\).
If \(S(x,y)\), eliminate the supplied finite disjunction
\eqref{eq:elim-global-cover}.  In each branch the corresponding
intersection theorem puts \((x,y)\) in the union.  Its length is bounded
by the sum of the intersection length bounds; its height maximum is a
bound for all target points.

No finite computation is said to decide the universal cover premise.
That premise is used with a proof, or retained as a hypothesis.  The
program only computes the lists justified by the displayed contracts.
\end{proof}

A pointwise auxiliary-existence assertion suffices when all its witnesses
belong to one explicitly bounded coefficient space.  This observation is
useful because no choice function or globally selected cover is needed.

For \(e,H\ge1\), let \(\mathcal B(e,H)\) be the finite list of nonzero
rational polynomials of total degree at most \(e\), with every reduced
coefficient numerator bounded in magnitude by \(H\) and every positive
denominator bounded by \(H\).  Zero coefficients are represented by
\(0/1\).  There are
\[
 K=\frac{(e+1)(e+2)}2
\]
monomials of degree at most \(e\).  Listing all numerator-denominator
pairs and then all coefficient vectors gives at most
\[
 B(e,H)=\bigl((2H+1)H\bigr)^K
 \tag{E10}\label{eq:elim-box-size}
\]
raw vectors; normalization and removal of zero polynomials can only
decrease that number.  Nonzero constants may also be removed, since they
cannot vanish at a target point.  A radix-vector scan implements this
enumeration by one fixed primitive recursive definition even when \(e\)
is an input.

\begin{theorem}[A bounded auxiliary envelope gives positive finiteness]
\label{thm:elim-envelope}
Assume
\[
 \forall x,y\in\Q\quad S(x,y)\ \longrightarrow\
 \exists G\in\mathcal B(e,H)\,
       \bigl(\Good(F,G)\ \wedge\ G(x,y)=0\bigr).
 \tag{E11}\label{eq:elim-pointwise-envelope}
\]
There is a fixed primitive recursive algorithm, depending on the fixed
predicate \(A\), which takes \(F,e,H\) and produces a complete list of
target points and a height or code bound.  A coarse bound on the number
of points is
\[
 d e\, B(e,H).
 \tag{E12}\label{eq:elim-envelope-card}
\]
In particular, the conclusion does not require a selector taking a point
to its auxiliary polynomial.
\end{theorem}

\begin{proof}
Enumerate the entire finite coefficient box.  Retain every \(G\) for
which the good-resultant test succeeds, and concatenate all corresponding
lists \(\operatorname{Intersect}(F,G)\).  Filter the concatenation by
\(A\) and remove duplicates.  Every retained point satisfies \(S\).

Conversely, given \(S(x,y)\), eliminate the bounded existential in
\eqref{eq:elim-pointwise-envelope}.  Its witness occurs in the enumerated
box, passes the test, and has \(G(x,y)=0\).  The complete intersection
list for that witness therefore contains \((x,y)\).  Hence the union
contains it, and the true predicate \(A(x,y)\) retains it.  This proves
completeness without selecting witnesses as a function of points.

There are at most \(B(e,H)\) candidates and each successful candidate
has degree at most \(e\), so its list has at most \(de\) entries.
This proves \eqref{eq:elim-envelope-card}.  Taking a maximum over the
computed complete list gives the height or code bound.  The bound on the
coefficient box is explicit, and every subsequent list-producing call is
primitive recursive; their composition is primitive recursive as well.
\end{proof}

The degree \emph{and coefficient} bounds are doing different jobs.  The
degree controls intersection cardinality.  The coefficient bound makes
the space of all possible witnesses a finite searchable list, which
produces a height bound and complete enumeration.

For comparison, suppose only that every finite list of target points can
be covered by at most \(M\) auxiliary equations of degree at most \(e\),
each with a nonzero elimination certificate.  The auxiliaries may depend
on the whole finite list and their coefficients may be unbounded.
Assign every list entry to its first covering equation.  Each group has
at most \(de\) points by Theorem~\ref{thm:elim-bezout}, hence every
distinct target list has length at most \(Mde\).  This is a valid finite
cardinality argument.  It supplies neither a global cover to use in
Corollary~\ref{cor:elim-global-cover} nor a coefficient envelope to use
in Theorem~\ref{thm:elim-envelope}.  In particular, changing auxiliary
equations along larger and larger lists does not give a stopping rule for
enumerating all target points.  The missing uniformity has been located,
not inferred away.

\subsection{Exact examples and what the construction repairs}

Take
\[
 F=X^2+Y^2-5,\qquad G=XY-2.
\]
The first passing shear has \(c=1\), giving
\[
 F_c=U^2-2UV+2V^2-5,\qquad G_c=UV-V^2-2,
 \qquad R(U)=(U^2-9)^2.
\]
For \(u=3\) the retained ordinates are \(1,2\); for \(u=-3\) they
are \(-1,-2\).  The inverse shear returns precisely
\[
 (1,2),\ (2,1),\ (-1,-2),\ (-2,-1).
\]
There are four points, attaining \(de=4\), although only two first
coordinates after the canonical shear.  Thus fibre multiplicity is
necessary for the sharp count in this example.

For \(F=XY-2\) and \(G=X+Y-3\), the candidate \(c=0\) fails the
leading-coefficient test for \(F\), and \(c=1\) fails it for \(G\).
The scan chooses \(c=2\).  Here
\[
 F_c=UV-2V^2-2,\qquad G_c=U-V-3,
 \qquad R(U)=-(U-4)(U-5)
\]
with our row convention.  The complete list is \((1,2),(2,1)\),
attaining \(de=2\).  This illustrates a presentation in which the
bounded shear is genuinely needed.

If \(F=Y-X\) and \(G=2Y-2X\), the resultant is zero and the algorithm
returns unsupported.  Returning an empty list would be false: the
intersection contains every rational point on that line.  The code also
handles elimination roots at zero by stripping the initial polynomial
power before the rational-root scan.

Relative to the audited packets, the construction adds three concrete
interfaces: a bounded conversion from arbitrary plane coordinates to
constant leading coefficients, a complete rational-point producer for
every successful intersection certificate, and a finite-envelope route
from bounded auxiliary existence to positive finiteness.  It does not
provide the quantitative approximation theorem or coefficient bounds
needed to apply that route to all curves in the general Siegel or
Faltings statements.  It also does not provide actual rational
Mordell--Weil quotient representatives or effective exceptional-zero
estimates for prime-distribution arguments.

\subsection{Quantifier-free induction audit and Lean completion}

All claims above have a direct first-order arithmetic form over natural
codes.  A displayed bounded conjunction means the corresponding
primitive recursive Boolean fold equals one.  Validity guards on rational
and polynomial codes prevent division by zero; invalid input branches
receive fixed defaults.  The following are the main proof obligations.

\begin{enumerate}
\item \textbf{Polynomial operations.}  Implement coefficient arrays,
finite evaluation, multiplication, shear, and synthetic division.  Prove
the division identity by induction on the remaining coefficient count.
The invariant is a coefficientwise equality of coded polynomials, with
finite support and degree bounds compiled into Boolean tests.
\item \textbf{Determinants.}  Enumerate length-\(N\) vectors over
\(\{0,\ldots,N-1\}\), retain vectors with distinct entries, and sum
their signed products.  The loop bound \(N^N\) is primitive recursive.
Prove specialization, the adjugate identity, determinant multiplication,
and column factor extraction by finite expansion and reindexing.  The
degree audit records the sum of selected column indices; it needs no
unbounded algebraic assertion.
\item \textbf{Root scans.}  Compute denominator products, first and last
nonzero indices, gcds, and the rectangles
\eqref{eq:elim-root-box}.  The scan invariant says that every stored
entry passes the exact equation test and that every passing candidate in
the processed prefix is stored.  Both assertions are bounded Boolean
predicates.  Instantiate the separately proved divisibility bound after
the scan to obtain completeness for an arbitrary rational root.
\item \textbf{Intersection lists.}  The nested-loop invariant expresses
membership in the accumulated union of the processed finite root
fibres.  Completeness for an arbitrary point follows by two instances of
the root-list theorem and Lemma~\ref{lem:elim-root}.  The sharp length
bound additionally needs Vandermonde and the finite grouped-factor
invariant \(R=P_jQ_j\), with all quotient coefficients stored.
\item \textbf{Coefficient envelopes.}  The coefficient-box enumeration
has an explicit radix bound.  The union invariant identifies stored
points with passing intersections in the processed prefix.  The
unbounded cover premise is held in the proof context and instantiated
only after this induction.  It is never an induction formula.
\end{enumerate}

In Lean, the native specification may use rational coefficient lists and
finite arrays.  The essential definitions are
\texttt{shear}, \texttt{canonicalShear}, \texttt{sylvester},
\texttt{elimination}, \texttt{good}, \texttt{rationalRoots},
\texttt{intersect}, \texttt{auxiliaryBox}, and
\texttt{envelopeList}.  Define
\[
\begin{split}
 \operatorname{Common}(F,G,p)&:
       F(p_1,p_2)=0\wedge G(p_1,p_2)=0,\\
 \operatorname{Complete}(L,P)&:
       \forall p\,(p\in L\leftrightarrow P(p)).
\end{split}
\]
The principal native theorem targets are precisely the statements of
Lemmas~\ref{lem:elim-shear}--\ref{lem:elim-rational-roots} and
Theorems~\ref{thm:elim-complete}, \ref{thm:elim-envelope}.  The
complete-list theorem can be finished before the fibre-multiplicity
module: it uses the common-point root lemma, not the sharp cardinality
bound.  Thus that longer determinant argument need not block the positive
enumeration result.

The object-language phase uses the same fixed arithmetic programs,
their compiled Boolean predicates, and the displayed quantifier-free
invariants to construct actual \texttt{Prf PRAAx} derivations.  A native
Lean theorem about rational arrays is not yet such a derivation, and
\texttt{pra\_iff\_native} is applied only after that object proof has
been built.  The accompanying exact Python code checks the concrete
examples and exceptional branches; it is a reference implementation, not
a proof of the universal statements or an iPRA proof checker.
